\documentclass[class=NCU_thesis, crop=false]{standalone}
\begin{document}

\chapter{摘要}

根據 2026 年初衛生福利部統計資料，台灣登記在案之聽覺機能障礙者已超過十四萬人，而台灣手語（Taiwan Sign Language, TSL）為其日常溝通的重要媒介。然而，現有網路手語資源多以單詞查詢或固定句型影片為主，當使用者欲表達完整中文句子時，常需逐詞查找並自行組合對應手語影片，導致資訊零散且缺乏連續性。為提升手語數位內容之可用性，本研究建置一套由中文文字輸入至虛擬人手語動畫輸出的模組化系統，將中文至手語詞彙序列轉換、手語動作資料檢索、動作參數串接與虛擬人渲染整合於同一流程中。

在系統設計上，本研究將手語動作資料建置與使用者動畫生成分為兩類流程。手語動作資料建置流程使用 EHM-Tracker 將既有手語影片轉換為包含人體姿態、手部動作與臉部資訊之參數化追蹤資料，並建立可由 gloss 查詢的詞彙索引；使用者動畫生成流程則以前端中文至自然手語詞彙序列轉換模組（Sign Language Translation, SLT）產生 gloss sequence，依序查詢並檢查對應之詞彙追蹤資料。系統接著透過 motion envelope 分析擷取各詞彙之有效動作區間，視設定對詞彙參數進行時間平滑，並於相鄰詞彙片段之間加入過渡幀，以降低多詞彙串接時的姿態跳動與動作不連續問題。最後，系統將串接後的連續動作參數輸入 GUAVA avatar rendering 模組，產生具人物外觀一致性之虛擬人手語動畫。

實驗部分，詞彙層級平滑實驗顯示，10 個樣本之 SMPL-X body pose 與 MANO hand pose normalized jerk cost 平均分別降低 99.90\% 與 99.87\%，表示高頻動作變化明顯減少；然而，hand-only PCK@0.10 平均由 0.9476 降至 0.8868，顯示平滑度與手部逐幀對齊精度之間仍需取得平衡。本研究亦比較 EHM-Tracker 於 100-step 與 1000-step 最佳化設定下之追蹤結果。結果顯示，1000-step 在雙手清楚可見樣本中通常能改善 reprojection error 與 PCK，但在單手低可見度或手部離開畫面情境下，較長最佳化可能造成誤追蹤風險。因此，本研究採用較保守之 100-step 設定建立手語動作資料庫。此外，本研究透過使用者問卷評估 25 支系統生成影片，共回收 53 份有效問卷；四項指標平均分數介於 3.80 至 4.01 之間，其中準確度與手勢清晰度表現較佳。問卷結果顯示，本系統已具備中文至虛擬人手語動畫生成之可行性，但在手部細節、臉部表情與句子層級自然度上仍有改善空間。整體而言，本研究證明結合 SLT、既有手語影片資料、EHM-Tracker 與 GUAVA 之模組化流程，可作為後續聽障友善溝通、手語學習與手語數位內容生成之基礎。

\vspace{2em}
\noindent \textbf{關鍵字：} \keywordsZh{} % Set keywords in config.tex
\end{document}
